\documentclass[11pt,letterpaper]{article}
\usepackage[margin=1in]{geometry}
\usepackage{amsmath,amssymb,amsthm,booktabs,graphicx}
\usepackage[T1]{fontenc}
\usepackage{lmodern}
\usepackage[compact]{titlesec}
\usepackage[colorlinks=true,linkcolor=black,urlcolor=blue]{hyperref}
\theoremstyle{definition}
\newtheorem{definition}{Definition}
\newcommand{\E}{\mathbb{E}}
\newcommand{\Var}{\operatorname{Var}}
\newcommand{\Cov}{\operatorname{Cov}}
\newcommand{\1}{\mathbf{1}}
\newcommand{\frc}{\operatorname{frac}}
\setlength{\parskip}{4pt}
\setlength{\parindent}{0pt}

\begin{document}
\begin{center}
{\Large\bfseries Rotation rounding: a stateless replacement for stochastic rounding}\\[4pt]
{\small Method note, 10 November 2026}
\end{center}

\textbf{In one paragraph.} When weights are trained directly on a low-bit grid, every update has to be rounded back
onto the grid, and the rounding errors pile up in the weights. Stochastic rounding (SR) rounds up with a fresh random
threshold each step: it is unbiased, but its errors are independent, so they never cancel and the accumulated error
grows like a random walk. Rotation rounding keeps SR's per-step behaviour exactly, but takes its thresholds from a
golden-ratio sequence, so a round-up is soon followed by a round-down. The errors largely cancel, the optimizer sees
less rounding noise, and it costs no extra memory. In our 2-bit training runs it beats SR at every state budget, by
up to 0.11 bits per byte.

\section{Setting}
Each weight lives on a grid with spacing $q$. An optimizer step moves it off the grid and it must be rounded back,
because only grid values are stored. Measured in grid units, the off-grid position is an integer plus a fraction
$p_t\in[0,1)$. A rounding rule picks a threshold $u_t\in[0,1)$ and rounds \emph{up} if $u_t<p_t$, otherwise down.
The rounding error at step $t$, in grid units, is
\begin{equation}
e_t=\1[u_t<p_t]-p_t ,
\end{equation}
i.e.\ $1-p_t$ if we rounded up and $-p_t$ if we rounded down. Since only the rounded value is kept, the errors add
up: after $N$ steps the weight sits $q\,S_N$ away from where the exact updates would have put it, where
\[
S_N=\textstyle\sum_{t<N}e_t .
\]
A rounding rule is just a choice of thresholds $u_t$, and we want it to keep $S_N$ small.

In our experiments the deployed weight has 2 bits (4 levels). A method with $b$ \emph{counter bits} also stores the
position between deployed levels on a grid $2^b$ times finer, and rounds onto that finer grid. Its training state is
$2+b$ bits per weight. With $b=0$, every update is rounded straight onto the 4 deployed levels.

\section{The two rules}
\begin{definition}[Stochastic rounding, SR]
Draw $u_t$ uniformly from $[0,1)$, independently at every step and for every weight.
\end{definition}
\begin{definition}[Rotation rounding]
Give each weight a random offset $u_0\in[0,1)$ and set
\[
u_t=\frc(u_0+t\varphi),\qquad \varphi=\tfrac{\sqrt5-1}{2}\approx0.618,
\]
where $\frc(x)=x-\lfloor x\rfloor$ is the fractional part. Each step the threshold moves forward by $0.618$ and wraps
around $[0,1)$. For example, $u_0=0.9$ gives $u_1=\frc(1.518)=0.518$, then $0.136$, $0.754$, $0.372$, \ldots
\end{definition}
\textbf{Storage.} Nothing per weight. The offsets $u_0$ are regenerated each step from a fixed random seed per weight
matrix; only the global step counter $t$ is stored. In code:
\begin{verbatim}
    g = torch.Generator().manual_seed(seed)   # same seed every step
    u = torch.frac(torch.rand(shape, generator=g) + (t*PHI) % 1)
    round_up = u < p
\end{verbatim}

\section{Why rotation is better}
\textbf{Step 1: each single step is exactly SR.} For a fixed $t$, $u_t$ is uniform on $[0,1)$, because $u_0$ is.
So rotation rounds up with probability $p_t$, just like SR. Each step is unbiased ($\E e_t=0$) and has the same error
variance $p_t(1-p_t)$. Rotation gives up nothing SR offers per step. \emph{This holds in any training run.}

\textbf{Step 2: SR's errors form a random walk.} SR draws a fresh threshold that is independent of everything before
it, including $p_t$, however the gradients reacted to earlier rounding. So each error has mean zero given the past,
the errors are uncorrelated, and their variances simply add:
\begin{equation}
\Var S_N=\textstyle\sum_{t<N}\E[p_t(1-p_t)] .
\end{equation}
This grows linearly in $N$, so the typical size of the accumulated error, $\sqrt{\Var S_N}$, grows like $\sqrt N$.
Nothing ever pulls it back. \emph{This also holds in any training run.}

\textbf{Step 3: rotation's errors cancel.} The golden-ratio sequence spreads its points very evenly over $[0,1)$: in
any stretch of consecutive steps, the fraction of thresholds below a level $p$ is close to $p$. So if the weight keeps
getting the same fraction $p$, rotation rounds up in almost exactly a fraction $p$ of the steps, and the errors
cancel. Figure~\ref{fig:sim}a shows this. $|S_N|$ stays bounded by $O(\log N)$, against SR's $\sqrt N$. This is the
same idea as error diffusion (sigma--delta modulation), but without storing the error.

\begin{figure}[h]
\centering
\includegraphics[width=\linewidth]{rotation_fig.pdf}
\caption{Simulated rounding errors (2000 weights, no training). (a) Constant fraction $p=0.3$: SR's sums wander off,
rotation's stay near zero. (b) Fractions $p_t$ drawn at random each step: rotation's error still grows like $\sqrt t$,
but $\sqrt{1-\rho}=\sqrt{1/2}$ times as fast, as Eq.~\eqref{eq:main} predicts (dashed lines).}
\label{fig:sim}
\end{figure}

In training, $p_t$ changes from step to step. Suppose the $p_t$ are random, drawn independently of the thresholds,
and write $F(u)=P(p>u)$. Then
\begin{equation}
\Var S_N = N\,(1-\rho)\,\E[p(1-p)] + O(\log^2 N),\qquad
\rho=\frac{\Var_{u\sim U[0,1)}F(u)}{\E[p(1-p)]}=\frac{\E\min(p,p')-(\E p)^2}{\E[p(1-p)]}\in[0,1],
\label{eq:main}
\end{equation}
where $p,p'$ are two independent draws. Compared with SR's $N\,\E[p(1-p)]$, rotation removes a fraction $\rho$ of
the random walk.\footnote{Proof sketch: split $e_t=m_t+(F(u_t)-\E p)$, where $m_t$ is the part due to the random
$p_t$; it has mean zero and is independent across steps. Summed over $t$, the second term is the error of
approximating $\int F$ by an average of $F$ over the golden-ratio points. Because $F$ is monotone and $\varphi$ has
bounded continued-fraction coefficients, the Koksma--Hlawka inequality bounds it by $O(\log N)$.} Intuitively, $\rho$
is the share of the error that is decided by the threshold rather than by the randomness of the update itself; only
that share can be cancelled.
\begin{itemize}\setlength{\itemsep}{0pt}
\item Same fraction every step: $\rho=1$; all of the random walk is removed (Figure~\ref{fig:sim}a).
\item Fraction uniform on $[0,1)$ (large steps spanning many grid cells): $\rho=\tfrac12$ (Figure~\ref{fig:sim}b).
\item Small noisy steps around a grid point, typical of training near convergence: $\rho\approx0.3$--$0.4$.
\end{itemize}
The irrational step $\varphi$ matters. With $\varphi=0$ (a fixed random threshold per weight, ``frozen dither''), the
same threshold repeats every step and the error grows linearly into a bias.

\section{Why the accumulated error is what matters for the weights}
Near a minimum with curvature $h$ and learning rate $\eta$, the distance from the optimum, $d_t=x_t-x^*$, follows
\[
d_{t+1}=a\,d_t+q\,e_t,\qquad a=1-\eta h,
\]
ignoring gradient noise. Unrolling, $d_T=q\sum_{k\ge0}a^k e_{T-k}$: the weight error is a decaying average of past
rounding errors, with a memory of about $1/(\eta h)$ steps (hundreds, typically). Its variance is
\[
\Var d=q^2\sum_{j,k}a^{j+k}\Cov(e_{T-j},e_{T-k})\;\approx\;\frac{q^2}{2\eta h}\sum_{k=-\infty}^{\infty}\Cov(e_0,e_k)
=\frac{q^2}{2\eta h}\,\lim_{N\to\infty}\frac{\Var S_N}{N},
\]
where the approximation needs the error correlations to die out well within the memory window. So the weights feel
the long-run growth rate of the error sum, not the per-step error:
\begin{itemize}\setlength{\itemsep}{0pt}
\item \textbf{SR}: the errors are uncorrelated, so only the $k=0$ term survives: $\E[p(1-p)]$.
\item \textbf{Rotation}: same $k=0$ term, but a round-up now makes a round-down likely within the next few steps, so
the terms $k\ne0$ are negative. They sum to $-\rho\,\E[p(1-p)]$, leaving $(1-\rho)\,\E[p(1-p)]$.
\end{itemize}
In the frequency domain, the sum of covariances is the error's power spectrum at zero frequency, while the per-step
variance is its total power. Both rules have the same total power. SR spreads it evenly over all frequencies (white
noise). Rotation moves it to high frequencies: fast up--down patterns that cancel within a few steps. The optimizer
acts as a low-pass filter with cutoff near $\eta h$, so it averages the high-frequency part away and only the
zero-frequency part reaches the weights. In the simulation of Figure~\ref{fig:sim}b with $\eta h=0.01$, rotation's
weight-error variance is $0.505$ of SR's, against the predicted $1-\rho=0.5$.

\section{What is proven and what is assumed}
Steps 1 and 2 are exact in any training run. Step 3 assumes that the fractions $p_t$ do not line up with the
threshold sequence. In training, $p_t$ depends on earlier rounding through the loss, so this is not automatic: an
adversary who knew $u_0$ could pick $p_t$ just above $u_t$ every step and force a round-up each time, making the error
grow linearly. The gradient, however, depends on $u_0$ only indirectly, through the weights that past rounding
produced, and nothing in it steers $p_t$ toward the next threshold. We therefore expect step 3 to hold in practice.
It is supported by a one-weight closed-loop simulation with feedback (rotation lowers the excess loss from rounding by
about 25\% at 0--2 counter bits) and by the training runs below, where rotation is never worse than SR. It is not proven.

\section{Results}
Two-bit weights trained from scratch, no momentum, same normalized step for both rules, learning rate tuned per
method, mean of 3 seeds. The metric is validation bits per byte (lower is better). The \emph{gain} is SR minus
rotation.
\begin{center}\small
\begin{tabular}{@{}llcccccc@{}}
\toprule
& counter bits $b$ & 0 & 1 & 2 & 4 & 6 & 8 \\
& state (bits/weight) & 2 & 3 & 4 & 6 & 8 & 10 \\
\midrule
2.65M model & SR & 2.570 & 2.383 & 2.228 & 1.976 & 1.938 & 1.918 \\
WikiText-2, 2.5k steps & rotation & 2.519 & 2.332 & 2.168 & 1.959 & 1.934 & 1.906 \\
& gain & 0.051 & 0.051 & 0.060 & 0.017 & 0.004 & 0.012 \\
\midrule
10.6M model & SR & 2.391 & -- & 2.055 & 1.721 & 1.636 & -- \\
WikiText-103, 8k steps & rotation & 2.334 & -- & 1.941 & 1.689 & 1.635 & -- \\
& gain & 0.057 & -- & 0.114 & 0.032 & 0.001 & -- \\
\bottomrule
\end{tabular}
\end{center}
On a teammate's format with learned per-group codebooks instead of a uniform grid (GPEF, 4 bits/weight), rotation also
beats SR by 0.05, so the effect does not depend on the grid being uniform. On that format, however, GPEF's own
deterministic choice of the nearest level is better than both (10.6M: 1.889, against 2.261 for rotation and 2.316
for SR). There the learned fp32 codebook levels already absorb the small updates, so the comparison that matters is
rotation against SR, not against nearest-level rounding.

\textbf{Where the gain is largest.} Rotation removes rounding noise, not gradient noise, so its gain should be roughly
\[
\text{gain}\;\approx\;\rho\times(\text{share of the weight error that comes from rounding}).
\]
The two factors move in opposite directions as counter bits are added:
\begin{itemize}\setlength{\itemsep}{0pt}
\item \emph{The rounding share falls.} For steps smaller than a grid cell, rounding variance scales like
$q\cdot|\text{step}|$, so it halves with each extra bit, and gradient noise takes over. By $b=6$ the gain is gone.
\item \emph{$\rho$ rises.} At $b=0$ a step is a small fraction of a grid level, so $p_t$ sits near 0 or 1 and
$\rho\approx0.3$. Each extra bit makes the same step twice as large in grid units, which spreads $p_t$ out and pushes
$\rho$ toward $\tfrac12$ and above.
\end{itemize}
Their product peaks at an intermediate budget, which matches the gain being largest around 0--2 counter bits. Two
caveats. On the small model the gains at $b=0,1,2$ are nearly equal; only the 10M $b=2$ gain stands out, and its
rotation runs vary by $\pm0.034$ across seeds, so its size is uncertain. And this explanation is an interpretation of
the theory; we have not measured $\rho$ inside training.

\textbf{Bottom line.} Rotation rounding is a drop-in replacement for stochastic rounding: identical per-step
statistics, no extra memory, provably better error cancellation when updates do not track the thresholds, and better
accuracy in every training comparison we have run, most of all at 2--4 bits of training state per weight.

\end{document}
